
\section{Limitations of CORD-19}\fbeg{}\begin{frame}{\ft{Limitations of CORD-19}}

\begin{sftxt}[.87\textwidth]

According to AI2, roughly 30\% of CORD-19's full-text documents were 
shared in structured formats such as JATS-XML, but the majority of 
publications were only available as PDF files.  AI2 therefore 
had to extract text from those PDFs in order to perform conversion 
to S2ORC JSON.  This is an intrinsically error-prone process 
due to anomalies within internal PDF text representation, 
such as hyphenation, ligatures, spurious space characters 
in some font styles, headers/footers and marginal content that 
needs to be isolated from the main text, etc.  Furthermore, 
technical content such as chemical formulae or linguistics 
annotations can be difficult to read properly without 
working off of structured markup, such as XML.  Tramscription 
errors limited the extent to which inter-document connections 
and conceptual links can be properly recognized.  Also, 
S2ORC only models documents at the paragraph level --- not 
individual sentences --- which can limit the application of 
NLP-related text mining techniques.  \textcolor{white}{AI2 recognized these 
limitations and even printed a \q{call to action} encouraging 
publishers to develop more rigorous text encoding and metadata standards.}

\vspace{-2em}AI2 recognized these 
limitations and even printed a \q{call to action} encouraging 
publishers to develop more rigorous text encoding and metadata standards.
   

\end{ftxt}

\end{frame}

